\chapter{Results}
\label{ch:results}
% Numerical tables are generated. Each generated file identifies its evaluated snapshot.
% The main RQ1--RQ3 tables use the final snapshot with attribution off and on.
% Development versions appear only in Appendix A (app:repairs).

This chapter reports what the harness of Section~\ref{sec:protocol} measured on the corpus of
Section~\ref{sec:corpus}, one research question at a time. Section~\ref{sec:res-baseline}
answers RQ1, Section~\ref{sec:res-defects} RQ2, Section~\ref{sec:res-after} RQ3 and
Section~\ref{sec:res-mutation} RQ4. Section~\ref{sec:res-round2} asks what the losses were and what the cause of each says about the claim (RQ5). Section~\ref{sec:deadlinks} reports the dead links the method found as a
by-product.

The numerical tables are generated from the harness result files. The summaries compiled by hand
are marked as such. The main must-conversion, loss, before-and-after and graph-composition tables
describe the final extractor snapshot. The baseline is that code without scope attribution, and the
attributed version is the same code with the pass enabled. The development versions of uigraph
and the corrections between them are listed in Appendix~\ref{app:repairs}.

Two conditions hold for every figure below. The source trees are extracted without their dependencies installed, and the corpus was selected with the tool in view. Section~\ref{sec:limitations} discusses both. The second makes the figures likely to be optimistic for Angular applications in general.

\section{The two levels of loss (RQ1)}
\label{sec:res-baseline}

RQ1 is answered at two levels. The first is whether the graph holds the screens the application
declares. The second is how much of the application's navigation becomes a deterministic
transition between them (the ratios are defined in Section~\ref{sec:metrics}). The first is a
precondition of the second, because a transition cannot start or end at a screen the graph does
not have. It is not a numerical bound on the second: one ratio counts screens and the other
counts constructs, so a missing screen costs only the constructs that name it.

Table~\ref{tab:conversion} gives the second measure on the final snapshot. The constructs that
earn must number \CensusBaseMustSites{} in the whole corpus without attribution and
\CensusMustSites{} with it. They lie in \CensusMustApps{} applications with attribution, and the
other \CensusZeroMustApps{} have none. Edge counts must be read per tier and not summed as
distinct constructs, because one unresolved site can produce many over-approximated may edges.

Why the tier is small is measured directly. A measurement switch of the final version records, at
every may witness, each condition that held it at may. It records them at the extractor's own decision points. With
the switch off the output is unchanged, which was checked on all six applications
(Section~\ref{sec:reproducibility}). Table~\ref{tab:mayreasons} groups the reasons of the
\MayOnlyConstructs{} constructs that yield may edges and no must edge. \MayAppAny{} of them sit behind
at least one decision of the application itself that the extractor reads: a guard or resolver, a
branch, a callback or a template condition. Only \MayToolOnly{} are held at may by the extractor's
own uncertainty alone, so a better reader could lift at most those to must. The small must tier is
therefore mostly a property of how these applications navigate, not of the extractor. Three limits
apply.
\begin{itemize}
\item A guard that always lets the navigation through still counts as a guard: the census says the
  application decides, not that the navigation is conditional at run time.
\item Mifos holds \MayOnlyConstructsMifos{} of the \MayOnlyConstructs{} constructs, so every share is
  mostly its own.
\item The families overlap. Only the \emph{sole} column counts single causes.
\end{itemize}

\begin{table}[t]\footnotesize
\caption{Why the constructs that yield only may edges are may, by family of reason. \emph{Any}
counts the constructs with at least one reason of the family, \emph{Sole} those whose every may
witness has that one reason only. The application's own decisions are the async, branch, template
and guard families; reader, attribution and target are the extractor's uncertainty; same screen
is neither.}
\label{tab:mayreasons}
\TabMayReasons
\end{table}

\begin{table}[t]\small
\caption{Must conversion per application, baseline. Must conversion is
the fraction of in-scope constructs whose own witness earns must; the edge columns
count edges. \emph{Nodes} counts every node of the baseline graph: the screens of
Table~\ref{tab:coverage} together with its unplaced origins and dynamic targets
(Table~\ref{tab:census} counts the attributed graphs, which add anchors and more of the latter). The baseline and attributed graphs are extracted from the same final
snapshot on identical corpus inputs.}
\label{tab:conversion}
\TabConversion
\end{table}

Table~\ref{tab:coverage} gives the first measure. It compares the nodes the extractor mints with
the screens the independent route parser finds declared. Against that denominator node
extraction is complete: on all six applications every declared screen has a node.

Section~\ref{sec:res-round2} gives the causes of the node-extraction defects the final version
reads, and Appendix~\ref{app:readers} the handful of addresses the two readers still disagree
about.

Two qualifications apply, because the figure invites a reading that neither supports.

\begin{itemize}
\item The denominator is the screens a \emph{static reader} can enumerate. Against the addresses
the running framework mounts, the same node set is \ProbeNodes{} of \ProbeScreens{}
(Section~\ref{sec:res-probe}). Almost all of the difference is one construct in one application
that the extractor declares unread.
\item The two static readers reached their agreement partly by being corrected against each other.
Their agreement is therefore weaker evidence than if they had stayed independent. The runtime
method is the external check for that reason, and it finds no node at an address where the
framework mounts nothing (Section~\ref{sec:res-unmounted} gives its provenance).
\end{itemize}

Completeness here is a claim about the screens the readers can \emph{see}. The blind-spot column marks where it is incomplete but does not bound it. One route entry neither reader can resolve can hide many screens (Section~\ref{sec:res-probe}). Such an entry is outside the denominator. The applications that have such entries say so in their own soundiness notes. Refusing to read
more of an application would raise this column, not lower it. It is printed beside the coverage
figure for that reason.

\begin{table}[t]\footnotesize
\caption{Node extraction against the independent route parser, in both directions: declared
screens that have a matching node, and nodes for which the independent route parser holds no declared
screen. The blind-spot column counts route entries the independent route parser could not follow, which marks where either direction is incomplete. A blind spot is a route the independent route parser
knows is there and cannot address: a lazy boundary whose route array it could not find, a path
expression it could not resolve to a literal, a route whose path is computed by a custom matcher, or a
route source spread into a table from a call it does not follow; the last kind now accounts for most of
the column. The node set is the same with scope attribution off and on, since attribution mints no
nodes.}
\label{tab:coverage}
\TabCoverage
\end{table}

\section{Which failure is worth an analysis (RQ2)}
\label{sec:res-defects}

Table~\ref{tab:defects} lists the defect classes and the evidence for each, separating the named
evaluation methods from the other observations. Four of the five accepted classes rest on one
named method each. Attribution has three observations on top of it, which gives four lines of
evidence. Guards, named outlets and constant paths have one observation each, which gives two. The
lazy-boundary class rests on three methods. The attribution class leans most on observations, and
it is the class the analysis of Chapter~\ref{ch:analysis} addresses. Its observations are not weak
in themselves, but they are not independent instruments.

The five classes were found during development, and each has at least two confirmations.
The quantitative breakdown revealed a sixth candidate: a node that exists but is not linked to the
component that renders it. It is confirmed only by the comparison between the independent route parser and the
graph. It is therefore reported as a candidate until a second method confirms it. It is empty on the
final version.

\begin{table}[tbp]\footnotesize
\caption{Defect classes and the evidence for each. A class is admitted on two
lines of evidence, at least one from a named evaluation method; the methods
column counts only those. This summary is compiled by hand; it is not harness output.}
\label{tab:defects}
\begin{tabularx}{\textwidth}{@{}p{0.27\textwidth}Yrr@{}}
\toprule
Class & Evidence & Methods & All \\
\midrule
Attribution: navigation outside route components is invisible & mutants excluded on the baseline graph; multi-mount edge absent; manual audit; an application with nodes and no edges on a development version (k8s dashboard; Appendix~\ref{app:repairs}) & 1 & 4 \\
Lazy module boundary not followed & independent route parser; test-source cross-check; mutation & 3 & 3 \\
Guards unreliable in both directions & red-team panels and a source survey of the guard factory and the call-form \texttt{canActivate}
(Section~\ref{sec:res-losses}) & 1 & 2 \\
Named router outlets not implemented & test-source cross-check; source survey & 1 & 2 \\
Constant-valued route paths not resolved & independent route parser; node-count collapse & 1 & 2 \\
\addlinespace
\emph{Candidate:} node not linked to its component & independent route parser against graph & 1 & 1 \\
\bottomrule
\end{tabularx}
\end{table}

Table~\ref{tab:loss} and Figure~\ref{fig:loss} assign every navigation construct to one outcome
by the rules of Table~\ref{tab:loss-rules}. Of the \RoundTwoInScopeAll{} constructs in the corpus,
they show how many convert onto each tier and why the rest do not. A construct counts on the must
tier only when its own witness earns that modality. The rows of Table~\ref{tab:loss-rules} are
loss categories, not the defect classes of Table~\ref{tab:defects}. Guards and named outlets are
measured outside this assignment. The categories that match no admitted defect class were not put through
the admission rule. Guard defects are measured separately (Section~\ref{sec:metrics}).

\begin{table}[t]\small
\caption{Outcome of every navigation construct, baseline. Columns: realworld,
alfresco, ghostfolio, k8s dashboard, bitwarden, mifos. All rows are measured on the
final snapshot with attribution disabled.}
\label{tab:loss}
\TabLoss
\end{table}
\FigLoss

Under the order in which the loss rules test a construct, attribution is the largest loss category in every application. Of the
\RoundTwoInScopeAll{} constructs, the two attribution categories hold \AttrLost{} between them,
and no other loss category holds a tenth as many. The constructs whose destination the extractor records as unresolved are counted apart. Some are run-time values and some are references the extractor declines to follow. They come to about a third as many.
Most of the attribution categories come from components that several screens render, not from
services.

This ranking depends on the order in which the loss rules test a construct. Ownership is tested
before the target, so a construct outside a route component is filed as attribution whatever its
target. Of the attribution constructs, \AttrNonLiteral{} have a target that the loss rules'
literal test does not accept. If the target were tested first, the target category would be
larger than attribution on \AttrOrderFlipWhere{}.

Constant route paths, lazy module boundaries and nodes with no link to the component that renders
them are empty in Table~\ref{tab:loss}, because the final extractor reads their causes
(Section~\ref{sec:res-losses} explains all three once). The ranking justifies spending an analysis on
attribution: once the extractor reads what the applications declare, the navigation it still
cannot place is mostly navigation with no single owning screen, under the order caveat above.

The categories also interact, in both directions. A construct in a shared component used only inside
screens with no node is counted as attribution, although attribution alone cannot recover it. A
construct whose own route component has no node is counted as a node defect, although attribution could
reach it through a shell or subtree scope. A residual node-missing category still holds a few constructs
(Table~\ref{tab:loss}), and attribution converts none of them (Table~\ref{tab:recovery-source}). That it
could is a possibility the rules allow and this instrument does not show. The \AttrLost{} attribution
constructs are therefore neither a floor nor a ceiling on what scope attribution recovers
(Section~\ref{sec:res-after}).

Table~\ref{tab:oracles} reports the witness check on the final attributed tier and the
test-source cross-check. Section~\ref{sec:res-round2} treats the must tier separately.

\begin{table}[t]\small
\caption{Results of the witness check, on the final attributed tier, and of the test-source
cross-check. Mutation is
reported separately in Section~\ref{sec:res-mutation}, on the attributed graph. This
summary is written by hand from the harness output rather than generated.}
\label{tab:oracles}
\begin{tabularx}{\textwidth}{@{}p{0.26\textwidth}Yr@{}}
\toprule
Method & Result & $n$ \\
\midrule
Witness check, final attributed must tier & \MustWitnessViolations{} detected
failures; \MustWitnessUnread{} edges have a destination the checker cannot read.
Unread edges are not passes, and target agreement does not establish
Claim~\ref{claim:must}. & \MustWitnessEdges{} \\
Test-source cross-check, alfresco & \EtoERecallAlfresco{} URLs the specifications name have a node; one
negative test excluded. & \EtoERowsAlfresco{} \\
\bottomrule
\end{tabularx}
\end{table}

\section{Scope attribution (RQ3)}
\label{sec:res-after}

The analysis under evaluation is the one Chapter~\ref{ch:analysis} states: scopes in Section~\ref{sec:scope}, the rules that compute them in Section~\ref{sec:attribution} and the anchored representation in Section~\ref{sec:anchored}.

Three quantities
decide the answer: how much navigation the pass recovers, whether it costs the deterministic tier
anything, and what it still cannot place. The third shows where the analysis as implemented stops.

Table~\ref{tab:before-after} and Figure~\ref{fig:before-after} compare the extractor with and
without scope attribution on identical inputs. The comparison is made construct by construct. A
construct is \emph{recovered} if it was lost to attribution in the baseline and is witnessed by an
edge after attribution.

\begin{table}[t]\footnotesize
\caption{Before and after scope attribution on the final evaluated snapshot.
Must conversion in per cent and its change in percentage points; must-tier edges whose destination
the witness check read and passed; and constructs assigned to attribution in the
baseline that recover onto must or may. A construct's tier is the modality its own
witness earns. The remaining must edges are unread destinations, all \MustWitnessUnread{} of them
on k8sdash, which are neither passes nor failures; no edge failed. Constructs converted from other baseline categories are in Table~\ref{tab:recovery-source}.}
\label{tab:before-after}
\TabBeforeAfter
\end{table}
\FigBeforeAfter

Attribution recovers \AttrRecovered{} of the \AttrLost{} constructs the loss rules assign to it.
This is \AttrRecoveredShare{} of them, which this thesis calls the attribution recovery. Most of them are on mifos. Without mifos the share is \AttrRecoveredShareNoMifos{}, and \AttrUnderHalfApps{} of the six applications recover under half. It takes the graphs from \AttrEdgesBefore{} edges to
\AttrEdgesAfter{}.

Measured by RQ1's must conversion, the gain is far smaller. Of the recovered
constructs, \AttrRecoveredMust{} earn must and \AttrRecoveredMay{} earn may, and \AttrAppsGained{}
of the six applications gain any must conversion at all. The must constructs rise from
\CensusBaseMustSites{} to \CensusMustSites{}. Two of the rise are those recoveries, both on
realworld. The other two are k8s dashboard constructs that the baseline already converted at may
and attribution raises to must. The recovery count does not include them.

The two are different measures of recovery, and both are reported. A gain in represented
navigation need not be a gain in deterministic transitions when targets or execution conditions
remain uncertain. Counted in edges, the must tier grows from \AttrMustBefore{} to \AttrMustAfter{}.
The edges added are the links of two application shells, drawn once for each screen the shell
renders (Section~\ref{sec:res-round2}).

Following each construct by its identifier across the two runs finds \AttrRegressed{}
that leave the converted categories, so no must conversion of the baseline was lost. This measures recovery
of represented constructs and says little about must-tier precision. The evidence for that tier,
and the edges outside it, are treated in Section~\ref{sec:res-soundness}.

The anchored representation was checked on the same graphs. Claim~\ref{claim:expansion} requires
the expansion of the anchored edge set to equal the materialised edge set. The anchored edges are
therefore materialised independently. Each scope-anchored edge is fanned over the screens its
anchor stands for, with its own deduplication. The reason is that the library's own materialisation function is
a loop over the query and could only disagree about which nodes it iterates. The two agree on every
application, and Table~\ref{tab:expansion} is not produced otherwise. The table is also
the only measurement of what the representation saves. On \ExpansionWidest{} the materialised edge
set is \ExpansionWidestRatio{} times the size of what is stored.

The check exercises the member walk and the deduplication key. The definition of a scope's membership
is shared by both paths and is not tested by it. The claim is therefore supported on \ExpansionApps{}
applications and on fixtures, not proved. On \ExpansionLibraryLarger{} of them the library's own
function returns a larger set than either, which is the iteration difference and not a disagreement
about the edge set.

\begin{table}[t]\small
\caption{The anchored edge set against its expansion. \emph{Anchored} is what the graph
stores, \emph{expanded} what the per-screen query yields, and \emph{independent} a second
materialisation written for this check. Claim~\ref{claim:expansion} requires the last two to
be equal; the check refuses to emit this table if they differ on any application.}
\label{tab:expansion}
\TabExpansion
\end{table}

Table~\ref{tab:recovery-source} gives the recovery by the loss category each construct held in the final
baseline and by the modality its own witness earns. Of the recoveries from the attribution categories,
\AttrRecoveredMust{} earn must and \AttrRecoveredMay{} earn may. The loss categories are not
independent causes: a shared component can be assigned to attribution even when a missing target
node prevents its recovery.

\begin{table}[t]\small
\caption{Where recovered constructs came from on the final snapshot. Rows are their
baseline loss categories; columns give the modality each construct's own witness earns
after attribution. Constructs are followed by identity, not matched by source
location.}
\label{tab:recovery-source}
\TabRecoverySource
\end{table}

Table~\ref{tab:remaining} classifies each construct that stays lost.

\begin{table}[t]\small
\caption{Why attribution losses remain after attribution.}
\label{tab:remaining}
\TabRemaining
\end{table}

The two tables close the attribution categories between them. \AttrRecovered{} recovered plus \AttrRemaining{}
remaining is the \AttrLost{} the rules assigned. The script that produces the tables asserts this before it
writes either table. All but \AttrRemainingUnexplained{} of the remaining constructs carry a reason. Those
\AttrRemainingUnexplained{} have a scope and a destination the graph holds, and the diagnosis names
no reason they stay lost. The reasons are limitations of the implemented analysis and limits of the
source population.

\begin{itemize}
\item Some components have no recovered route or template path to a screen. This is an empty
computed scope, not proof that the application never renders them.
\item Some sites lie beyond the member-level call graph's supported receiver and provider forms.
\item Other sites have a scope but a dynamic destination, or a target for which the graph has no
node.
\end{itemize}

\subsection{Checking the scopes against the running interface}
\label{sec:res-scope}

Everything above measures what attribution recovers. Whether the \emph{set of screens} it computes
for a site is the set that site fires from is a different question, and the contribution depends on
it. The rendered-interface check of Section~\ref{sec:metrics} answers one direction. Table~\ref{tab:scope} retains the recorded
observations of an earlier graph. The browser experiment was not repeated and its predictions were
not recomputed for the final call graph, so this section documents what that earlier graph missed.

\begin{table}[t]\small
\caption{Recorded scope predictions of an earlier graph against the rendered
interface, on the \ScopeApps{} applications the harness could serve. These predictions
were not recomputed for the final call graph. \emph{Judged} counts screens that rendered themselves rather than
redirecting to a sign-in page; the rest cannot be judged from what rendered, because a
redirected screen shows another screen's links. \emph{Not predicted} is the only column that
counts against the analysis; the other direction is reported in the text, because it can be
judged only where the graph names a destination exactly.}
\label{tab:scope}
\TabScope
\end{table}

A scope claim can be too small or too large. Too small means the application renders a link the
graph does not attribute to that screen, which the rendered interface decides. Too large means the
graph attributes a site to a screen the application never fires it from, which the interface decides
only where the graph claims a destination exactly. The two are reported separately.

\paragraph{Scopes too small.} Of the \ScopePairs{} (screen, target pattern) pairs that the running
applications render on the judged screens, \ScopeConfirmed{} are predicted by the computed scope
(\ScopeConfirmedShare{}). Behind these predictions stand \ScopeAnchors{} raw anchors. On
realworld, which renders most of its screens without a backend, every rendered pair on every
judged screen is predicted. The pairs that are not predicted are all bitwarden's.

The unit of the check is the screen and the destination, not the site. The graph's edges leaving a
screen are pooled and set against the links the screen renders. A confirmed pair therefore shows
that the screen's outgoing destinations cover what it renders, not that each site behind them was
attributed to that screen. An unpredicted pair is not affected by the pooling, so the
\ScopeUnpredicted{} are direct evidence of scopes that are too small.

The evidence is limited in size. It covers \ScopeJudged{} screens of \ScopeScreens{}, because
\ScopeUnjudged{} redirect to a sign-in page when there is no backend behind them. Only
\ScopeApps{} of the six applications are measured. ghostfolio is left out: it builds a separate
bundle per locale, each with a base path of its own, and the harness serves a single directory at the
server root. The check settles one direction of a question, on part of the corpus.

The \ScopeUnpredicted{} unpredicted links share one cause, and all four are in bitwarden. Its
\path{/login} screen renders links to \path{/hint} and \path{/signup}, and two further screens
render a link to \path{/signup}. The graph holds those edges, with the right witness
(\path{libs/auth/src/angular/login/login.component.html}), but anchors them at an origin node rather
than at a screen. The component that renders them lives in a workspace library outside the
application directory, and attribution computed the empty scope for it. The failure is a scope the
analysis declined to compute, which the running application shows to be non-empty. These sites are
the \emph{unplaced} class of Table~\ref{tab:round2acct}.

\paragraph{Scopes too large.} This direction does not support a single figure. \ScopeUnrendered{}
predictions run this way: the graph holds the edge and the page showed no link. Two of the three
causes mixed in that count are not defects. The harness signs in to nothing, so realworld's header,
wrapped in \texttt{@if (authState === 'authenticated')}, hides \emph{New Article}, \emph{Settings}
and the profile link (interface state). A tag page navigated to with a placeholder segment lists no
articles (data). The third cause is target over-approximation, a property of the destination and not
of the scope. realworld's \texttt{[routerLink]="['/profile', article.author.username]"} resolves to
\path{/profile/:username} and to \path{/profile/:username/favorites}. The deeper route is
recorded as a may edge at half confidence, which is the graph stating that it does not claim the
address.

Only predictions that name a destination exactly can be refuted by a page not linking there. There
are \ScopeExact{} of them across the judged screens, against \ScopeApprox{} over-approximated.
\ScopeOverAttributed{} of the \ScopeExact{} were not rendered. A denominator of \ScopeExact{}
supports no rate. The over-attribution direction was therefore not refuted on this
corpus, because the graph almost never claims an exact destination at a screen this check can
observe. The measure is reported as several numbers, not a precision figure, so that the two
directions stay separate.

\subsection{What the attribution discriminates}
\label{sec:res-width}

The rendered-interface check cannot say whether the scopes are too large, so this subsection asks
how much narrower the computed attribution is than the widest one there is. Assigning every
navigation site to every screen recovers every construct and needs no analysis. It is useless, because the
per-screen query would return the whole site list on every screen. The ratio between the two
measures what the analysis contributes (Table~\ref{tab:scopebaseline}).

\begin{table}[t]\small
\caption{The computed attribution against assigning every site to every screen. A site is a
source construct a witness cites; an edge anchored at a scope is counted over the screens that
scope stands for, since counting it at its anchor would understate the analysis's breadth.
\emph{Mean scope} is the attributed pairs per site.}
\label{tab:scopebaseline}
\TabScopeBaseline
\end{table}

The analysis is between \BaselineFactorMin{} and \BaselineFactorMax{} times narrower than that
baseline, and the low end of the range is the figure that qualifies RQ3. On \BaselineBluntest{}
the average site is attributed to \BaselineBluntestScope{} of \BaselineBluntestScreens{} screens,
which is \BaselineBluntestShare{} of the application. The per-screen query there excludes about two
screens in five and admits the rest. The widest margin is \BaselineSharpest{}, where
\BaselineSharpestAttributed{} attributions are computed where the baseline would assert
\BaselineSharpestTrivial{}. Across the corpus the mean scope runs from \BaselineMeanScopeMin{} to
\BaselineMeanScopeMax{} screens.

Two things follow, and the second is the less favourable to the analysis.

\begin{enumerate}
\item The trivial baseline is not a competitor for the must tier. An unguarded site in a shared
component attributed to every screen would yield a must edge from screens that never render it, so
Claim~\ref{claim:must} would fail by construction. That none of the readable must edges attribution
added fails the witness check is meaningful only where the attribution is narrow. Of the
\the\numexpr\AttrMustAfter-\AttrMustBefore\relax{} must edges it added, \MustWitnessUnread{} are
unread, all on \BaselineBluntest{}, where the scopes are widest.
\item On part of this corpus the attribution is not very narrow. A scope that names three fifths of
an application is a weak statement about where a site fires, and the attribution recovery RQ3 reports is bought
with scopes of that width. Nothing in the evaluation shows those wide scopes to be wrong, because the
rendered-interface check could not test the direction in which they would be wrong. Nothing shows
them to be right either. The recovered share of Table~\ref{tab:round2acct} should be read together
with the widths in this table.
\end{enumerate}

\subsection{What each mechanism contributes}
\label{sec:res-ablation}

The width says how far the attribution is from doing nothing, not whether the mechanisms that
produce it earn their place. Table~\ref{tab:ablation} extracts the corpus at \AblationLevels{}
nested levels, each withholding every mechanism above it. The withholding covers the stand-ins of
Section~\ref{sec:implementation} as well as the implemented relations. Otherwise the lowest
level would credit direct ownership with sites that only the stand-ins place.

\begin{table}[t]\footnotesize
\caption{The corpus extracted at each attribution level. \emph{Sites} counts the source
locations a witness cites that the level places on at least one screen. One construct can have
more than one, and a location can lie outside the application directory, so site counts are not
comparable with the construct counts of Table~\ref{tab:loss}. \emph{pairs} the (site, screen)
attributions it asserts. The levels are nested, and the table is produced only if
each one's pairs are a subset of the next one's. \emph{Screens} counts the screens at an
address a link can name, so the catch-all screens of alfresco and mifos are left out of it and of
the pairs, and the column falls short of the census's total by those screens.}
\label{tab:ablation}
\TabAblation
\end{table}

The levels are not evenly spaced.

\begin{itemize}
\item Direct ownership places \AblationOwnedSites{} sites, \AblationOwnedSiteShare{} of what the
full analysis places. Containment adds \AblationNestingAddedSites{} sites, the largest gain in
sites of any step, and brings the total to \AblationNestingSites{}.
\item The shell rule is the largest step in (site, screen) pairs and in reference recall, and this
thesis does not add it. \emph{Rooted} places \AblationRootedSites{} sites,
\AblationRootedSiteShare{} of the full figure, and carries the pair count from
\AblationNestingPairs{} to \AblationRootedPairs{}.
\item Everything above that level raises the site count from \AblationRootedSites{} to
\AblationFullSites{}, that is \AblationAboveRootedSites{} sites or \AblationAboveRootedShare{} of
the corpus total. Outlet descent (which an ancestor-route closure in earlier work also covers) adds
\AblationOutletAddedSites{} of them. The full level adds \AblationFullAddedSites{},
\AblationFullAddedShare{} of the corpus total. That is the member-level call graph that places a
service's navigation, its fallback join and the stand-ins for guards and effects together, all of
which only that level carries.
\end{itemize}

The call graph's own share is smaller. Against the coarser class-level rule it replaced, it places \AblationCallGraphNetSites{} more sites,
\AblationCallGraphNetShare{} of the corpus total, and \AblationCallGraphNetPairs{} more pairs.
That rule remains as the fallback for a member no caller reaches. The two
differ in pairs on \AblationClassChangedApps{} of the six applications: \AblationClassOnlyPairs{} pairs
only the class join asserts and \AblationFullOnlyPairs{} only the call graph does.

The shell rule is not a mechanism this thesis introduced, so the comparison with earlier work uses
two of the levels. Rima et al.'s specification also closes a route's
components over those of its ancestor routes, which is what \emph{outlet} adds. Both levels
reimplement those mechanisms inside this extractor and are not runs of either tool. The comparison
measures what the added mechanisms place on the same inputs, and does not show that the tools
themselves would place fewer. Against the nearest of them, \emph{outlet}, this thesis adds the
\emph{full} step alone: \AblationFullAddedSites{} sites, \AblationFullAddedShare{} of the corpus
total, every one outside a component, in a service, guard or effect. Rima et al.'s following of a
handler into same-class helpers and into a \texttt{subscribe} callback is already inside
\emph{outlet} here. The reason is that the extractor places every navigation call in a component's class at
that component's screens whichever member holds it.

Outlet descent contributes mostly width. It places \AblationOutletSites{} sites against rooted's
\AblationRootedSites{}, the same count on \AblationOutletSameSitesApps{} of the six applications,
while raising the pair count from \AblationRootedPairs{} to \AblationOutletPairs{}. A table of
counts cannot say whether that width is correct, which the next subsection addresses.

\subsection{Agreement with a structural reference}
\label{sec:res-groundtruth}

For \GTApp{} the screens each navigation site fires from were read off the sources by hand
(Section~\ref{sec:metrics}): \GTSites{} sites and \GTPairs{} (site, screen) pairs. Scoring every
level against that reading gives a comparison no other method here can make, because it has both
directions rather than one. Table~\ref{tab:groundtruth} reports it.

Two limits qualify the reading beyond those of Section~\ref{sec:metrics}.
\begin{itemize}
\item It is not independent of the graph. Its first version was completed
with nine header sites the scorer reported from the graph.
\item It is dominated by the shell. \GTShellSites{} of its \GTSites{} sites are placed on every screen and supply \GTShellPairs{} of the
\GTPairs{} pairs. Every level from \emph{rooted} upward places these pairs trivially.
\end{itemize} Over the other
\the\numexpr\GTPairs-\GTShellPairs\relax{} pairs, recall is \GTRootedRecallNonShell{} for
\emph{rooted}, \GTOutletRecallNonShell{} for \emph{outlet} and \GTFullRecallNonShell{} for
\emph{full}.

\begin{table}[t]\small
\caption{Each attribution level against the hand-read structural reference for \GTApp{}.
\emph{Missing} is a pair the reference has and the level does not; \emph{excess} is a pair
the level has and the reference does not. Precision and recall are agreement with that
reference, not with what the running application renders.}
\label{tab:groundtruth}
\TabGroundTruth
\end{table}

The table is produced only if each level's pairs are a subset of the next, and the
full level asserts no pair outside the reference. No lower level can then assert one either. The
levels differ only in recall. The absence of excess pairs at the lower levels follows from the
nesting and is not a separate result. Direct ownership recovers \GTOwnedRecall{} and containment
\GTNestingRecall{}, both far short. The shell rule, the level that stands for the prior work's
reach, recovers \GTRootedRecall{} at \GTRootedPrecision{} precision and leaves \GTRootedMissing{}
pairs unfound. Precision is measured against a reference whose header sites were added from the
graph's output, so it is weakest as evidence where the shell dominates.

Everything above \emph{rooted} reduces \GTRootedMissing{} missing pairs to \GTFullMissing{} and
asserts no pair outside the reference. The
last row, the class-level join that the call graph refines, asserts \GTClassExcess{} pairs outside the
reference, all at one site, the \texttt{navigate} inside \texttt{UserService.logout}. This takes
precision to \GTClassPrecision{}. The class join attributes a navigating service to every screen
whose components inject it. Ten components inject this one, the header among them. The shell
renders the header on every screen, so the site was attributed to every screen although its only
caller is a button on the settings screen. The call graph follows the call from that
button and places the site on the settings screen alone.

The clean row has one limit. The structural reference guided two repairs of the analysis. One is the call
graph, written after the reference exposed the class-level defect. The other is the rule that places an
empty-path child on the screen above it (Section~\ref{sec:approximations}). The \GTFullPrecision{}
precision and the recall of \emph{full} are scored on the same reference, so they are consistent with
the two repairs and test neither independently.

On realworld most of the recovery is produced by mechanisms earlier work already had:
\GTRootedRecall{} of the \GTPairs{} pairs are reached by the shell rule and \GTOutletRecall{} once the
enclosing routes are added. The \emph{full} step adds \GTOutletMissing{} pair and takes recall to
\GTFullRecall{}. realworld has one navigation inside a service, so it barely exercises the step only
this thesis adds.

\paragraph{A second application.} A second reference was therefore built for alfresco, where
navigation inside services occurs. Its protocol, the same rules as realworld's, was committed before
the reference was built or any level was scored (\path{ground-truth/ALFRESCO-PROTOCOL.txt}). The AI
assistant built it from the sources alone, without the graph. It places \GTAlfSites{} sites and leaves
\GTAlfUndetermined{} undetermined with the reason: on the viewer's named outlet, in components no route
places, behind actions declared in the extension configuration, or behind actions no source
dispatches. Those are not scored. Four lines the scorer then reported as placed by the graph and absent
from the reference were read afterwards under the same rules, given only the line numbers and not the
graph's placements. None is a navigation site, and they were added as excluded, where the protocol
provided only for adding a missed site; no figure depends on it, since unlisted lines are not scored.

\GTAlfShellSites{} of the sites, \GTAlfShellPairs{} of the \GTAlfPairs{} pairs, sit in the shell. The
application takes the shell from a library outside its sources and fills it through its extension registry. No level
reads either, so every level misses all but at most one of these pairs. Overall recall is therefore
\GTAlfOutletRecall{} at the levels of earlier work and \GTAlfFullRecall{} for the full analysis. Over the other pairs, where attribution decides, the levels that stand for earlier work recall
\GTAlfOutletRecallNonShell{}, the class-level join \GTAlfClassRecallNonShell{} and the full analysis
\GTAlfFullRecallNonShell{}. The gain is the service join. The two sites of \texttt{navigateBack} are
placed exactly on the two search screens their only caller renders. The full analysis asserts
\GTAlfFullExcess{} pairs outside the reference, all at \GTAlfFullExcessSites{} site,
\texttt{navigateToSearch}. Its only caller is an effect triggered by a toolbar button declared in the
extension configuration, which the analysis does not read, so it places the site on every screen at
the may tier. Most of the pairs the full analysis adds over \emph{outlet} come from this placement. The
class join's \GTAlfClassExcess{} excess pairs are at the same site, on the two search screens, where it
does not fire. The full and class levels also place one undetermined site on every screen,
\texttt{displayPreview} in \path{viewer.effects.ts}, whose action no source dispatches. It is not
scored. Scored as firing from nowhere, it would add to the excess of both. That takes precision to \GTAlfFullPrecision{}, against \GTAlfOutletPrecision{} for the
lower levels and \GTAlfClassPrecision{} for the class join. Of the protocol's predictions, the first,
that the full level recalls more than \emph{outlet}, held. The second, that its precision is at least
the class join's, did not.

What the two references establish is narrower than soundness. Both apply the analysis's own relations,
and neither licenses the corpus-wide widths of Table~\ref{tab:scopebaseline} on an application of six
hundred screens.

\subsection{Agreement with the applications' own suites}
\label{sec:res-e2e}

The end-to-end check of Section~\ref{sec:metrics} puts part of the question to the running
applications, with a reference that does not pass through the analysis's relations. The second and third validity checks hold on both applications, and the first on mifos, whose uninstrumented run fails the same tests. For realworld only instrumented runs were recorded, none of them with a failure.

\paragraph{realworld.} The suite completes \EtoERealworldEnds{} navigations, which reduce to
\EtoERealworldPairs{} distinct pairs, and the graph contains all \EtoERealworldHit{} of them:
\EtoERealworldHitScreen{} through an edge from the source screen or from a scope that names it, and
\EtoERealworldHitTop{} only through the shell scope, which names every screen. This result does not
bear on what the levels above \emph{rooted} add, because the \emph{rooted} level, which stands for earlier work's
mechanism, already contains \EtoERealworldHitRooted{} of the pairs, that is all of them. Without
attribution the graph contains \EtoERealworldHitNoattr{}. The suite exercises \EtoERealworldPairs{}
of the \EtoERealworldGraphPairs{} pairs the expanded graph contains, out of
\EtoERealworldPossiblePairs{} possible ordered pairs of screens including self-pairs. The
agreement is therefore with the transitions the suite takes and not with the graph.

\paragraph{mifos.} The suite completes \EtoEMifosEnds{} navigations, \EtoEMifosPairs{} distinct pairs,
and the graph contains \EtoEMifosHit{} of them. This count applies the pre-registered rule: a
redirected navigation is compared on the address it asked for. \EtoEMifosRedirect{} pairs were
redirected, and \EtoEMifosHitRedirect{} of those are hits. The scorer's
departure from this rule during development is listed in Table~\ref{tab:repairs}. Without attribution the graph contains
\EtoEMifosHitNoattr{}. None of the matched pairs is a must edge, as predicted. \EtoEMifosHitRedirect{} hits, those of the redirect bucket, rest only on
a site whose target prefix was over-approximated to many routes. They therefore show that the target is
reachable in the graph and not that the site was resolved. No other pre-registered prediction is
reported here.

\paragraph{The mifos misses.} Of the \EtoEMifosMiss{} misses, \EtoEMifosMissRedirect{} fall in the
pre-registered redirect class and are not classified further. The other classes were defined after the
misses had been read, and are tested in the order redirect, base, dynamic, first match winning. Every
one of the \EtoEMifosMiss{} misses meets the dynamic test whatever the order. The dynamic count is
therefore only what is left after the others. The test fires on any unresolved site in the
source component or above it, so it does not show that the missed navigation is that site. The pre-registered
five-class scheme and its prediction that misses would be mostly a scope too narrow or a site not read
are not scored.

\begin{itemize}
\item \EtoEMifosMissDynamic{} misses are classed \emph{dynamic}. They are places the extractor does
not read, not destinations only the run time knows. Three sites, in the clients, groups and savings
views, open an action screen with a template literal inside a commands array, for example
\texttt{navigate([\textasciigrave{}actions/\$\{name\}\textasciigrave{}], \{relativeTo\})}. The
reader treats that element as a variable and records the target as unresolved
(Section~\ref{sec:approximations}), although the text is static up to the action name. In the loan
view, \texttt{relativeTo} is passed through an options variable the reader does not evaluate, so the
relative prefix is matched from the root and names no route.
\item \EtoEMifosMissBase{} misses are classed \emph{base}, and are losses of a declared approximation, not repaired. The client
view's template holds three relative links, to the personal data tab and to the forms for a new loan
and a new savings account. From the client screen the graph resolves them correctly. The template is
rendered with each child tab, though, and attributed to the subtree scope of that screen. Its links
are then resolved against the child's address, where they name nothing mounted. The edges lead to the
not-found screen, and the graph lacks the edge the suite takes from the general tab. Angular resolves
such a link against the route of the component that holds it. This is the base resolution
Section~\ref{sec:deadlinks} names as this analysis's gap.
\end{itemize}

None of the misses was repaired: no other application with a runnable suite remains to test a repair,
and a repair fitted to this check would be untested by it.

\subsection{What scope attribution adds}
\label{sec:res-earlier}

This subsection gathers the RQ3 evidence in one place. Table~\ref{tab:earlier} sets three graphs side
by side: without attribution, with the mechanism of earlier work, and with the full analysis. The
earlier tools were not run. The \emph{rooted} column reimplements their mechanism inside
uigraph, on uigraph's own node extraction (Section~\ref{sec:res-ablation}). Every figure is in-sample.

\begin{table}[t]\footnotesize
\caption{The graph without attribution, with the mechanism of earlier work (\emph{rooted}) and with
the full analysis. \emph{Suite transitions} are the distinct pairs of screens the application's own
end-to-end suite navigates between (Section~\ref{sec:res-e2e}); \emph{graph pairs} are those the
graph holds. Only the full column of the suite rows was preregistered. The no-attribution graphs and
realworld's rooted graph come from existing measurements; the mifos rooted graph was mapped for this
comparison after the suites had been scored. A dash marks a figure not measured. All cells are
generated figures.}
\label{tab:earlier}
\begin{tabularx}{\textwidth}{@{}Yrrr@{}}
\toprule
 & no attribution & rooted & full \\
\midrule
Suite transitions in the graph, realworld & \EtoERealworldHitNoattr{} of \EtoERealworldPairs{} & \EtoERealworldHitRooted{} of \EtoERealworldPairs{} & \EtoERealworldHit{} of \EtoERealworldPairs{} \\
Suite transitions in the graph, mifos & \EtoEMifosHitNoattr{} of \EtoEMifosPairs{} & \EtoEMifosHitRooted{} of \EtoEMifosPairs{} & \EtoEMifosHit{} of \EtoEMifosPairs{} \\
Graph pairs, realworld & \EtoERealworldGraphPairsNoattr{} & \EtoERealworldGraphPairsRooted{} & \EtoERealworldGraphPairs{} \\
Graph pairs, mifos & \EtoEMifosGraphPairsNoattr{} & \EtoEMifosGraphPairsRooted{} & \EtoEMifosGraphPairs{} \\
Sites placed on at least one screen, corpus & -- & \AblationRootedSites{} & \AblationFullSites{} \\
Structural-reference recall, \GTApp{} & -- & \GTRootedRecall{} & \GTFullRecall{} \\
Structural-reference recall, alfresco & -- & \GTAlfRootedRecall{} & \GTAlfFullRecall{} \\
Structural-reference recall off the shell, alfresco & -- & \GTAlfRootedRecallNonShell{} & \GTAlfFullRecallNonShell{} \\
Structural-reference precision, alfresco & -- & \GTAlfRootedPrecision{} & \GTAlfFullPrecision{} \\
\bottomrule
\end{tabularx}
\end{table}

Three things follow.

\begin{itemize}
\item \emph{Attribution puts more of the suites' transitions into the graph, in a wider graph.} Without
  it, the graph holds \EtoERealworldHitNoattr{} of the \EtoERealworldPairs{} pairs realworld's suite
  takes and \EtoEMifosHitNoattr{} of mifos's \EtoEMifosPairs{}. With it, it holds \EtoERealworldHit{}
  and \EtoEMifosHit{}. The pairs the graph holds grow from \EtoERealworldGraphPairsNoattr{} to
  \EtoERealworldGraphPairs{} and from \EtoEMifosGraphPairsNoattr{} to \EtoEMifosGraphPairs{}. The
  extractor was fixed before the suites' logs were read (Section~\ref{sec:limitations}).
\item \emph{The gain is already reached at the rooted level, the mechanism of earlier work.} Oshima and
  Kishi's stated restrictions admit none of these applications: all \EarlierWorkParamApps{} declare
  route parameters and all \EarlierWorkNestedApps{} nested route trees. Rima et al.\ state no such
  restriction, and their specification covers this level (Section~\ref{sec:dynamic-spa}). Neither tool
  was run, so how they read these applications is not shown. Inside uigraph the level runs on all six
  because uigraph's route extraction reads parameters and nested trees, with or without attribution.
\item \emph{The service join places navigation no earlier mechanism can place.} A navigation call
  inside an injected service belongs to no component, so the levels up to \emph{outlet} cannot place
  it on any screen. The corpus has \AttrServiceLost{} such constructs at the baseline, on every one of
  the six applications (Table~\ref{tab:loss}). The service join, which only this thesis adds,
  recovers \AttrServiceRecovered{} of them and places \AblationFullAddedSites{} sites in services,
  guards and effects. On \GTApp{} its member-level call graph is also more precise than a class-level
  join: the class join asserts \GTClassExcess{} pairs outside the reference (\GTClassPrecision{}), the
  call graph none (\GTFullPrecision{}). On alfresco, where such navigation occurs, it raises recall off
  the shell from \GTAlfOutletRecallNonShell{} at the levels of earlier work to \GTAlfFullRecallNonShell{}
  (overall \GTAlfOutletRecall{} to \GTAlfFullRecall{}). However, the precision order reverses:
  \GTAlfFullPrecision{} against \GTAlfClassPrecision{} for the class join
  (Section~\ref{sec:res-groundtruth}). With the enclosing routes, the levels
  above \emph{rooted} find the \GTRootedMissing{} reference pairs rooted misses
  (Section~\ref{sec:res-groundtruth}). On the two applications with suites, they add no transition the
  suites take.
\end{itemize}

Of the \AttrLost{} constructs the loss rules assign to attribution, it recovers \AttrRecovered{}
(\AttrRecoveredShare{}), \AttrRecoveredMay{} of them at the may tier. The may tier is mostly the
applications' own decisions, not the extractor's uncertainty (Table~\ref{tab:mayreasons}). The
attributed must tier has \MustWitnessViolations{} witness-check failures among \MustWitnessEdges{}
edges, with \MustWitnessUnread{} destinations unread, which are not passes
(Section~\ref{sec:res-soundness}).

The scopes are wide where a shell renders a link on every screen: the computed attribution is
\BaselineFactorMin{} to \BaselineFactorMax{} times narrower than assigning every site to every screen
(Section~\ref{sec:res-width}). Wide is not the same as too large. Where the scopes were read against a
reference, the full analysis asserts no pair outside it on \GTApp{} and \GTAlfFullExcess{} at one site on
alfresco, where it also places an undetermined site on every screen. Both references apply the analysis's own relations, and only two of the six applications
were read this way. Section~\ref{sec:limitations} discusses both as a limitation of the
work. A scope that is too large still yields edges whose witnesses
agree, so no witness-based method can find it. The suites test only the missing direction, and only on
the transitions they take. Attribution also widens what a source change can be scored against.
\MutBaseExcluded{} mutants cannot be scored without it, against \MutExcluded{} with it. On the
\MutPaired{} mutants both graphs can score, detection is \MutPairedRate{} with attribution and
\MutPairedBaseRate{} without. The \MutPairedDisagree{} mutant they disagree on comes from how the
harness matches a guarded screen, not from the graph (Section~\ref{sec:res-mutation}), so the
detection rate per scored change is essentially unchanged.

Table~\ref{tab:world} puts the comparison in terms of what can be done without this thesis and with
it, on this corpus.

\begin{table}[t]\footnotesize
\caption{What the corpus shows can be done without this thesis and with it. The dynamic row is this
thesis's own rendered-interface check run without signing in, not a crawler with credentials; the
earlier tools were not run; every figure is in-sample. All cells are generated figures.}
\label{tab:world}
\begin{tabularx}{\textwidth}{@{}YYY@{}}
\toprule
 & Without this thesis & With uigraph \\
\midrule
Static extraction of the corpus applications & Oshima and Kishi's restrictions admit \EarlierWorkAdmittedApps{} of \EarlierWorkApps{} & extracts all \EarlierWorkApps{} \\
Screens checked without signing in & a browser run renders \ScopeJudged{} of \ScopeScreens{}; the rest redirect to sign-in & every screen the static readers find declared has a node; nothing is run \\
Dead links found from the source & -- & \RoundTwoDeadLinkCandidates{} candidate sites in \RoundTwoDeadLinkWhere{}, checked against the runtime route oracle (Section~\ref{sec:deadlinks}) \\
Transitions on \GTApp{} & its own suite takes \EtoERealworldPairs{} pairs & the graph holds \EtoERealworldGraphPairs{}, an over-approximation \\
Response to a source change & -- & \MutFallibleDetected{} of \MutFallibleScored{} mutants, at constructs the graph represents \\
\bottomrule
\end{tabularx}
\end{table}

\section{Change sensitivity (RQ4)}
\label{sec:res-mutation}

Mutation testing requires the application to be extracted again once per mutant, so this part of
the evaluation was run on three of the six applications. They are realworld and alfresco, which
carry the earlier mutation figures, and bitwarden, whose routes carry many guards, guard factories
among them. Every figure in this section comes from the attributed graph with harness
\HarnessVersionMutation{}, a fixed seed and a quota of eight mutants per operator. The corrections
made to the mutation harness are listed in Table~\ref{tab:repairs}.

Table~\ref{tab:mutation-sample} shows how the sample was drawn. The harness found \MutCandidates{}
candidate mutations, sampled \MutSampled{} under the quota and scored \MutScored{}. \MutExcluded{}
were excluded, all of them guard removals. For \MutExcludedUnreached{} no edge of the unmutated graph enters any
screen the guarded route yields. \MutExcludedNoScreen{} guards a route that declares neither a
component nor children. An excluded mutant carries no information about sensitivity. The rule removes
exactly the cases in which a removal could not be observed. Therefore \MutGuardScored{} of the
\the\numexpr\MutGuardScored+\MutExcludedGuard\relax{} sampled guard mutants are scored. The guard
rate below is conditional on the graph entering the guarded screen. The exclusions are few on the
attributed graph but not on the graph without attribution. Of the \MutBaseSampled{} mutants sampled
against it, \MutBaseExcluded{} are excluded and \MutBaseScored{} remain. This is the same recovery gap seen
from the mutation side.

\begin{table}[t]\small
\caption{Mutation sample. Candidates found, mutants sampled under the per-operator
quota, and mutants excluded as confounded with a recovery gap.}
\label{tab:mutation-sample}
\TabMutationSample
\end{table}

Table~\ref{tab:mutation} gives the detection rate per operator, and Figure~\ref{fig:mutation} the
outcomes summed over the applications. The rename operator is set aside first, because its detection
is close to guaranteed by construction: a rename mutant is excluded unless the unmutated graph already
holds a node for the segment it renames, and a screen is identified by its full path
(Section~\ref{sec:model}). It tests node extraction and supplies \MutRenameScored{} of the
\MutScored{} scored mutants. The three other operators are detected in \MutFallibleDetected{} of
their \MutFallibleScored{} scored mutants: \MutLinkDetected{} of \MutLinkScored{} deleted
\texttt{routerLink} bindings, \MutRetargetDetected{} of \MutRetargetScored{} retargeted
\texttt{navigate} calls and \MutGuardDetected{} of \MutGuardScored{} removed guards. A deleted \texttt{routerLink} is sampled only on a line an edge of the unmutated graph cites, so it is missed only if the extractor misreads its own input. The other two operators are constrained in the same way. A retarget is generated only for a \texttt{navigate} call whose first element is a plain absolute string at a line the graph already cites, and it writes one declared route as a literal. A guard removal is scored only where an edge already carries that guard and enters the guarded screen. The samples are small, in single digits per operator and application, so the pattern
says more than the rate.

\begin{table}[t]\small
\caption{Detection rate per mutation operator on the attributed graph, as detected over
scored mutants. A dash means the operator found no candidate in that application.}
\label{tab:mutation}
\TabMutation
\end{table}
\FigMutation

Deleting a \texttt{routerLink} moved the witness set rather than the edge set in
\MutLinkWitnessOnly{} of its mutants, because a second identical link in the same template still
produces the edge. This is why detection is judged on witnesses as well as edges. The operator found
no candidates in alfresco, which binds \texttt{[routerLink]}. Against the graph without attribution
the same sample yields only \MutLinkBaseScored{} scored link mutants, of which \MutLinkBaseDetected{}
are detected, because the candidates in the application shells produce no baseline edge.

A retargeting mutant replaces the whole commands array of a \texttt{navigate} call with a declared
route that names no parameter. It is detected when the edited line witnesses an edge to the new
destination and none to a destination it reached before. The destination is drawn from the declared
routes alone, so both configurations patch each site alike (Appendix~\ref{app:repairs} lists the
operator's earlier artefacts).

Every guard removal the run can score is detected: \MutGuardDetected{} of \MutGuardScored{}, on all
three applications, each judged by whether some edge loses a guard the removed list names. A guard is
recorded on the edges that enter the guarded screen, so a removal can be observed only where the graph
holds that screen and enters it. All \MutExcludedGuard{} exclusions are that condition failing.
The harness matches a class by the file the routing module imports it from, because bitwarden declares
two classes named \texttt{AccountComponent}. Among the scored guards the empty path makes no
difference: \MutGuardEmptyPathDetected{} of the \MutGuardEmptyPathScored{} on a route whose own path
is the empty string are detected.

One limit of the same matching remains, and it shows in the run without attribution. realworld's
\texttt{editor} route has two children, at the empty path and at \texttt{:slug}, that load the same
component file, so the harness cannot tell their screens apart. Without attribution no edge enters
\path{/editor}, but one enters \path{/editor/:slug}, so the guard on the empty-path child is scored
and missed rather than excluded. That is the one paired disagreement, and it is a limit of the
harness, not of the graph.

The mutation run neither confirms nor contradicts the defect class \emph{guards unreliable in both directions}; Chapter~\ref{ch:discussion} gives the reading.

Of the \MutScored{} scored mutants, \MutDetected{} were detected, and \MutBaseDetected{} of
\MutBaseScored{} without attribution. On the \MutPaired{} mutants both graphs can score, the
figures are \MutPairedDetected{} and \MutPairedBaseDetected{}. In this sample the graph responded to
every scored change in the constructs it represents.

\section{What the losses were (RQ5)}
\label{sec:res-round2}

The defect classes of Section~\ref{sec:res-defects} are a diagnosis. RQ5 asks what lies behind each
class and what that says about the claim, and it is answered by cause analysis on the final version.
The cause trace of Section~\ref{sec:metrics} defines the three outcomes used below.

\subsection{What the losses were}
\label{sec:res-losses}

This section keeps two words apart: a \emph{class} is a row of Table~\ref{tab:defects}, and a
\emph{cause} is what the source review found behind a class.

\paragraph{Read by the final version.} The information was in the source, so the loss was the
tool's.

\begin{itemize}
\item \emph{Lazy module boundary} and \emph{constant-valued route paths}: causes in the extractor,
read by the final version, and empty in Table~\ref{tab:loss}.
\item \emph{Guards}: the two unread forms (a guard on a route a wrapper returns, a
\texttt{canActivate} element written as a call) are read, and the guard rule is the router's own
semantics (Section~\ref{sec:res-soundness}).
\item The candidate class \emph{node not linked to its component}: empty on the final version.
\end{itemize}

\paragraph{Still not read.} The final version does not read these causes.

\begin{itemize}
\item \emph{Attribution}: \AttrRecoveredShare{} of its constructs are recovered. The remainder is in
Table~\ref{tab:remaining}, including sites that no relation the prototype implements reaches. It mixes gaps of the tool with
components that no screen may render, and no measurement separates them.
\item \emph{Named router outlets}: a construct the model represents and the extractor does not
implement. Its navigations are excluded from the denominator (\RoundTwoExcludedOutlet{} sites) and
its screens left out of the runtime route oracle's mounted set (\ProbeOutletScreens{} route
entries).
\item ghostfolio's spread of computed routes is a construct the extractor could read and does not.
\ProbeMissedGhostfolio{} of the \ProbeMissed{} mounted addresses without a node sit behind it.
\item Three readable forms that the end-to-end check exposed: a template literal inside a commands
array, \texttt{relativeTo} passed through a variable, and the base resolution of a template shared
with child tabs (Section~\ref{sec:res-e2e}).
\end{itemize}

\paragraph{Beyond a static reader of a route table.} No reader of a route table can resolve these.

\begin{itemize}
\item ghostfolio's custom \texttt{UrlMatcher} computes a route's path in a function. No reader of a
route table can resolve it, and both tools report it as a blind spot.
\item mifos chooses between two route tables by a value injected at deployment, which the runtime
method of Section~\ref{sec:res-probe} finds and the sources leave undecided.
\end{itemize}

The causes the final version reads fall into three groups, traced here in the source.

\begin{sloppypar}
\paragraph{The route table was not reached.} This is the lazy-boundary cause. An NgModule loaded by
\texttt{loadChildren} keeps its routes in a separate \texttt{*-routing.module} that it imports.
A reader that stops at the module file finds nothing, and if the routing module is then scanned at
the project root, every screen in it is minted at its own relative path. Four further forms hide a
table.
\begin{itemize}
\item A route array held inline by \texttt{RouterModule.forChild([\ldots])}, whose elements may name
routes (\texttt{[ERROR\_ROUTE]}) rather than write them.
\item A loader that references its module directly
(\texttt{loadChildren: () => VaultModule}).
\item A barrel that re-exports the route array.
\item A wrapper
such as mifos's \texttt{Route.withShell([\ldots])}, a call whose array argument is skipped whole.
\end{itemize}
A lazy route may also name a module through a barrel that re-exports a feature module, which imports
the routing module that holds the table. A reader that follows one hop finds no table and places the
screens at the application root (\path{/breach-report}), where the running router mounts them deeper
(\path{/reports/breach-report}). Both static readers shared this cause, so their agreement could not
expose it and the runtime route oracle did. The final extractor follows the chain to the table,
bounded at three hops. At every hop it refuses a file the route itself sits in, a file another route
loads, and a file whose every table is registered at the application root. A related cause left nodes
with no link to the component that renders them. It is the candidate class and is empty on the final
version.
\end{sloppypar}

\paragraph{The path was named, not written.} This is the constant-path cause. In one application every
route path is a reference into a constant table, which collapses its graph to a single node when it
is not followed.
\texttt{path: internalRoutes.account.path} is a compile-time string, and so are an enum member and a
\texttt{\$localize} segment. Following the name needs the workspace's \texttt{paths} aliases and the destructuring through which a route
file reaches the table. The aliases matter because
the table commonly lives in a library outside the application directory.

\paragraph{The destination was named, not written.} A target may be a field, a dotted chain into
the route table, an element of a commands array, the return of a builder, or a value in a lookup
the component fills. Resolving these is what separates a link from a dead end. A field seeded with a
literal is mutable state, not a constant, and the reader does not follow it, because following it
would turn a run-time value into a must edge.

Node extraction is therefore not the bottleneck on the final version: every application is complete
against the readers' denominator. An analysis that decides where a navigation can fire needs a node
for the screen it fires from, and a construct whose source screen has no node is a loss attribution
cannot repair. 

\subsection{Construct accounting}
\label{sec:res-accounting}

Table~\ref{tab:round2acct} gives the construct accounting in four parts, because a single ratio cannot
separate three different things. A site is \emph{recovered} when it becomes a transition with a
destination. It is \emph{dynamic} when the extractor records the destination as unresolved: a value from the backend or a history entry, but also a field it treats as mutable although the route table initialises it, or a template literal inside a commands array (Section~\ref{sec:approximations}). The graph then records an edge into that screen's
\texttt{u\_}\,sink with modality unknown, rather than dropping or guessing. It is \emph{unplaced} when
the destination is readable but the screen it fires from is not, and the edge is anchored at an
\texttt{o\_}\,origin node that says so. What is left carries a recorded reason but no transition. The
check enforces that a reason is recorded, not that it is correct.

\RoundTwoExcluded{} sites leave the denominator, on three grounds decided from the construct's own text
before the graph is consulted. The set is therefore the same in both runs, and the figures are not produced
if the grounds do not sum to the total. \RoundTwoExcludedQueryOnly{} pass an empty commands
array and so change only the query of the current route, and \RoundTwoExcludedNoCommands{} names no
destination. Neither is a transition between screens. \RoundTwoExcludedOutlet{} navigate a secondary
outlet. These are transitions that the model represents as a second mount
(Section~\ref{sec:inputs}) and the extractor does not implement, so they are a declared recall
limitation. Excluding them raises alfresco's ratios. Both runs use the same ledger, so the two
columns of Table~\ref{tab:before-after} are ratios over the same \RoundTwoInScope{} sites. The
comparison refuses to print if any construct changes its ground between them.

\begin{table}[t]\small
\caption{Construct accounting after the loss list was worked. Percentages are of the
in-scope sites.}
\label{tab:round2acct}
\TabRoundTwoAccounting
\end{table}

Of \RoundTwoInScope{} in-scope sites, \RoundTwoRecovered{} are recovered, \RoundTwoDynamic{} are
accounted as dynamic and \RoundTwoUnplaced{} as unplaced, so \RoundTwoAccounted{} are accounted for in
all. k8s dashboard shows why recovered and dynamic
must be reported separately: a low recovered fraction can coexist with many sites whose destination
depends on runtime data. Its dynamic sites are spread over \RoundTwoDynamicFiles{} files,
\RoundTwoDynamicFilesOne{} of which contribute a single site, so the column reports a habit of the
codebase and not one place to fix. mifos supplies \RoundTwoMifosShare{} of the in-scope sites on its
own and pulls the corpus figures towards its own: without it, recovered is \RoundTwoNoMifosRecovered{}
and accounted \RoundTwoNoMifosAccounted{}. Changes to the route denominator are recorded in
Appendix~\ref{app:validation}.

The named-outlet restriction is enforced on the route side as well. A route that declares a named
\texttt{outlet} is reached as \texttt{(name:path)} and never at the primary address, so it contributes
no screen. Otherwise it could take the primary address from the sibling that really renders there, by
being declared first. The graph would then name the wrong component at a real address. bitwarden
declares such routes (a login and a registration screen each pair a primary component with a second
one on a named outlet at the same empty path), but the rule moves no count there. The restriction
reaches a named outlet's \emph{subtree} too. alfresco mounts its viewer as
\texttt{\{ outlet: 'viewer', children: [\ldots] \}}. The routes beneath it were minted at the primary
address, \path{/view/:nodeId} and \path{/search-libraries/view/:nodeId}, which Angular reaches only as
\texttt{(viewer:\ldots)}. Both tools read the subtree the same way, so the coverage figure could not
expose it, and the runtime route oracle of Section~\ref{sec:res-probe} was needed. Neither tool mints such
addresses.

\subsection{The final tier and the node set}

The attributed must tier has \MustWitnessViolations{} witness-check failures among \MustWitnessEdges{}
edges, with \MustWitnessUnread{} destinations unread; Appendix~\ref{sec:res-soundness} reports the
soundness property construct by construct. Across the corpus the compared graph declares
\ProbeInvented{} screens at addresses the framework mounts nothing at
(Appendix~\ref{sec:res-unmounted}).

\section{Dead links, found from the source}
\label{sec:deadlinks}

A site whose target the application mounts nothing at is a dead link. The method finds candidates
for such links at no extra cost. The candidates are found statically, and each is then checked
against the route configuration the runtime route oracle recorded before it is reported.

The clearest case is mifos. A transactions tab links to a journal entry with

\begin{quote}
\texttt{[routerLink]="['/', 'journal-entry',}\\
\texttt{\ \ \ \ \ \ \ \ \ \ 'view', 'S' + transaction.id]"}
\end{quote}

\begin{sloppypar}\noindent The application does declare that screen, at
\path{/accounting/journal-entry/view/:id}, inside the accounting module. The link omits the parent
prefix, so the URL it builds is one no route matches. mifos declares
\texttt{\{ path: '**', component: NotFoundComponent \}}, so the link does not fail invisibly. It
renders the not-found screen, which a user can mistake for missing data rather than a broken link.\par\end{sloppypar}

\begin{sloppypar}
The same shape repeats across the application.
\begin{itemize}
\item Links to \path{/loans-accounts/...} name a subtree
declared only under an owner (\texttt{/clients/:clientId/}, \texttt{/groups/:groupId/}).
\item Links to
\path{/savings-accounts/...} name one declared only under \texttt{/centers/:centerId/},
\texttt{/clients/:clientId/} and \texttt{/groups/:groupId/}.
\item Further calls name segments the application mounts under other names, such as \path{/users/create} (mounted as \texttt{appusers}) and \path{/system/scheduler-jobs} (\texttt{manage-jobs}).
\end{itemize} A test that never opens the tabs holding
these links would not find them, and a crawler would find them only if it reached those screens.
Reading the route configuration finds them without running anything.
\end{sloppypar}

The list is not one defect class, and the split matters for what each part is worth.

\begin{itemize}
\item \RoundTwoDeadLinkAbsolute{} of the candidates are an absolute link written against a path the
application does not mount, usually the name the server uses rather than the one the router does.
Those are the application's.
\item \RoundTwoDeadLinkRelative{} come from a \emph{relative} \texttt{routerLink}, such as
\texttt{[routerLink]="'edit'"} or \texttt{[routerLink]="['../../transactions']"}. Their target
depends on resolving the link against the screen it fires from, and getting that base wrong is this
analysis's gap before it is the application's defect.
\item \RoundTwoDeadLinkBaseUnknown{} are neither, because the result files do not record enough of
the construct to say which it is. They are counted and not classified, rather than assigned to the
half that would flatter the finding.
\end{itemize}

The last two groups are reported as candidates and weighted lower. Counted as places in the source
rather than as targets, the three groups are \RoundTwoDeadLinkAbsoluteLocations{},
\RoundTwoDeadLinkRelativeLocations{} and \RoundTwoDeadLinkBaseUnknownLocations{}. The difference is
one relative link in a shared template that resolves against three owning screens, so a single
mistake yields three candidates at one place in the source.

A further \RoundTwoOutsideDirSites{} sites are not published at all. The reason is the boundary of
the measurement, not the way the sites were read. They sit outside the directory the measurement is
scoped to, in \RoundTwoOutsideDirWhere{}, in the repositories of \RoundTwoOutsideDirApps{}
applications. Their files enter the analysis only because those repositories declare a
\texttt{paths} alias for them. Most are the source of a \emph{sibling front end}, a browser
extension built from the same monorepo. Its own router declares the same addresses a census would
have called dead.

Publishing them would report one application's valid links as another's defects. They are
counted here and left out of the census.

